\documentclass[11pt]{article}
\usepackage[margin=1in]{geometry}
\usepackage{amsmath,amssymb,amsthm,booktabs}
\usepackage[T1]{fontenc}
\usepackage[utf8]{inputenc}
\usepackage{lmodern}
\usepackage[hyphens]{url}
\usepackage[colorlinks=true,linkcolor=blue,citecolor=blue,urlcolor=blue]{hyperref}
\newtheorem{theorem}{Theorem}
\newtheorem{lemma}[theorem]{Lemma}
\newtheorem{proposition}[theorem]{Proposition}
\newcommand{\PAF}{\operatorname{PAF}}
\newcommand{\Z}{\mathbb Z}
\setlength{\emergencystretch}{2em}
\title{Nonexistence of Williamson matrices of order 65}
\author{P. Graham\\ \small\url{https://orcid.org/0009-0007-8493-1505}}
\date{October 22, 2026}
\hypersetup{pdfauthor={P. Graham},pdftitle={Nonexistence of Williamson matrices of order 65}}
\begin{document}
\maketitle
\begin{abstract}
Williamson matrices of order $n$ are four symmetric circulant sign matrices
whose squares sum to $4nI_n$. Their construction dates to Williamson's
1944 work on Hadamard matrices. We give a computer-assisted nonexistence
result at order $65$, described by Bright in 2019 as ``currently the smallest
order for which it is unknown'' whether Williamson matrices exist
\cite[Section 4.2]{BrightStatement}. Order $67$ remains open.
The search uses row sums, spectral inequalities, and compression to
reduce the problem, then exhaustively matches pairs of full rows.
Recorded case and pair totals are checked separately from the search.
The conclusion also depends on correct execution of the pair matching,
which has not been formally verified.
\end{abstract}

\section{Introduction}
A \emph{sign matrix} has entries in $\{1,-1\}$. A matrix is circulant when
each row is the cyclic shift of its predecessor. Four symmetric circulant
sign matrices $A,B,C,D$ of order $n$ are \emph{Williamson matrices} if
\begin{equation}\label{eq:matrix}
 A^2+B^2+C^2+D^2=4nI_n,
\end{equation}
where $I_n$ is the identity matrix. We establish the following
computer-assisted result.
\begin{theorem}\label{thm:main}
There are no four symmetric circulant sign matrices $A,B,C,D$ of order
$65$ satisfying $A^2+B^2+C^2+D^2=260I_{65}$.
\end{theorem}
This is a statement about symmetric circulant blocks, not about arbitrary
Hadamard matrices of order $260$ or the broader Williamson-type families.

\smallskip
\noindent Exhaustive computer search with a replayable CPU implementation,
receipts and two independent coverage and accounting checkers.

\noindent Computer-assisted; not yet reviewed by experts.
The proofs, the search code, the certificate checkers and the text were
produced by AI models (from OpenAI and Anthropic), working under the
author's direction.

\subsection{History and the new case}
We recall the existence question and explain how the present computation
extends earlier classifications. Williamson introduced the construction in 1944 \cite{Williamson}.
The odd-order Williamson conjecture asserts that these four blocks exist
for every odd positive integer $n$; Turyn recorded the more general
all-orders conjecture in 1972 \cite{Turyn}.
Baumert, Golomb and Hall's computer construction at block order $23$
yielded a Hadamard matrix of order $92$ \cite{BGH}.
Baumert and Hall subsequently enumerated the odd block orders through
$23$ \cite{BH}. Later searches by Sawade, Koukouvinos and Kounias,
and \DJ okovi\'c \cite{Djokovic1992,Djokovic1995} extended the odd-order classifications; the history,
including corrections to incomplete enumerations, is reviewed by
Bright, Kotsireas and Ganesh \cite[Section 1]{BKG}.
\DJ okovi\'c disproved the conjecture at order $35$ in 1993
\cite{Djokovic}. Holzmann, Kharaghani and Tayfeh-Rezaie completed the
odd-order enumeration through $59$, including nonexistence at $47$, $53$
and $59$ \cite{HKT}. Their work dates to a 2007 manuscript \cite{HKT}. Turyn's construction supplies orders $61$ and $63$.
Lang and Schneider also classified order $61$, as credited by
Bright, Kotsireas and Ganesh \cite[Section 1]{BKG}, in their work on
Turyn-type Williamson matrices \cite{LS}.
Bright, Kotsireas and Ganesh's MathCheck SAT and computer-algebra work, first posted
in 2018, enumerated all even orders through $70$ and the odd
multiples of three through $69$, including a new class at order $63$
\cite{BKG}. Bright's March 20, 2019 research statement describes $65$ as
``currently the smallest order for which it is unknown'' whether Williamson
matrices exist \cite[Section 4.2]{BrightStatement}; it also lists $67$ as
unenumerated. The present theorem answers the order-$65$ case negatively;
order $67$ remains open. The Wikipedia article on the Williamson conjecture
lists nonexistence at $35,47,53,59$; it does not identify $65$ as open.\footnote{\url{https://en.wikipedia.org/wiki/Williamson_conjecture}, consulted October 22, 2026.}
The open status of $65$ cited here rests on Bright's 2019 statement and
version 2 (September 6, 2026) of Kharaghani, Mohammadian, Tayfeh-Rezaie and Zaitsev's paper
\cite[Section 2.3]{KMTZ}.

There is substantial recent progress for a related family. Kharaghani,
Mohammadian, Tayfeh-Rezaie and Zaitsev's version 2 \cite{KMTZ}
constructs \emph{near-Williamson} matrices at order $65$: these satisfy
$AA^{\mathsf T}+BB^{\mathsf T}+CC^{\mathsf T}+DD^{\mathsf T}=260I_{65}$,
with $B,C,D$ symmetric, while $A$ need not be symmetric. The work was
first posted on May 9, 2026 as a three-author version 1 by Kharaghani,
Mohammadian and Tayfeh-Rezaie, titled \emph{A search for Hadamard matrices
of Williamson type}, with existence through odd order $63$.
The four-author version 2, dated September 6, 2026, establishes near-Williamson existence
at every odd order through $69$ and enumerates that broader family through
odd order $47$. Its order-$65$ example does not have symmetric $A$.
Our theorem excludes only the construction with four symmetric circulant blocks.
Section 2.3 of version 2 \cite{KMTZ} still includes $65$ among the Williamson
orders for which matrices have not been found; its order-$65$ search
constructs near-Williamson examples and does not exhaust the symmetric case.
The recent survey by Egan, Kotsireas, London and \'O Cath\'ain, first posted
on October 19, 2026, places these developments in the broader Hadamard
construction problem \cite{EKLO}; it records the known nonexistence orders
$35,47,53,59$, without reporting a settlement of the order-$65$ case.

Bennett, Bright, Colinot and Nayak's work, first posted January 29, 2026,
enumerates circulant Williamson-type matrices through order $21$, allowing
nonsymmetric blocks \cite{BBCN}. Their correspondence with quaternionic
perfect sequences gives an enumeration algorithm for that broader family.
This is relevant progress in Williamson-type enumeration; the present
order-$65$ search concerns symmetric blocks and uses neither their code
nor their quaternionic correspondence.

Zhang's April 2, 2026 report \cite[Section 7.4]{Zhang} describes an
AI-assisted attempt on the exact order-$65$ problem and expressly leaves
it unsolved. This earlier attempt is relevant background; none of its
code, data, or proposed reductions is used here.

The new contribution is completion of the order-$65$ problem with no
restriction on the four symmetric circulant blocks beyond their defining
equation. Compression and spectral restrictions are established
methods, developed in particular by \DJ okovi\'c and Kotsireas
\cite{DK} and used in the MathCheck SAT and computer-algebra enumeration
\cite{BKG}. The difficult step here is excluding \emph{every} lift of
a possible compression, meaning every choice of full rows having that
compression. Excluding the compressed objects alone does not settle the
question. A direct join of pairs of full rows, with conservative spectral bounds,
completes this step. This replaces the SAT-based lifting step of
Bright, Kotsireas and Ganesh's MathCheck system, whose published enumeration treats even
orders and orders divisible by three. Their enumeration does not include
$65$, which has neither of those properties.
Holzmann, Kharaghani and Tayfeh-Rezaie already used row and pair spectral
bounds \cite[Section 3]{HKT}. After choosing the first two blocks, they
used the product identity to eliminate the fourth and solved linearized
equations modulo $4$ for the third, checking the original equations
finally. Here compression, a multiplier restriction on one row, and
two-sided pair bounds instead precede a direct join of the two full-row
pairs. These are combinations of established tools, not a new general
construction principle. The measured cost below describes this
implementation; no controlled comparison with the earlier programs is
available, and no speedup over them is claimed.
The central limitation is explicit: checking the records is much cheaper than repeating the search, but does
not prove correctness of the filtering and matching operations.

Section~\ref{sec:overview} gives the argument's structure.
Sections~\ref{sec:rows}--\ref{sec:join} justify the finite reductions and
matching test. Section~\ref{sec:computation} states the computational
outcome, controls, replay instructions, and precise trust boundary.

\section{Proof overview}\label{sec:overview}
We first describe how the matrix equation becomes a finite enumeration.
A symmetric circulant matrix of odd order is determined by half of its
first row and its initial entry. Negating a whole matrix leaves its square
unchanged, so each initial entry may be fixed to $+1$. At order $65$ this
leaves a $32$-bit mask for each row. Row sums and Fourier eigenvalues give
necessary conditions that can be checked separately for each row.

Next, sum entries in each residue class modulo $13$. This compresses a
row of length $65$ to an integer vector of length $13$. Partition the
retained rows into classes according to their row sum and compression.
Every solution determines an admissible tuple: four classes satisfying
the exact compressed autocorrelation equations. A lift of that tuple is
a choice of one full row from each class. Common multiplication of the indices
by an invertible residue gives equivalent tuples, so only one tuple from
each orbit is needed. A search unit is one such representative. This is a cover of the full problem because every solution can be carried
to one of the representatives by the permitted symmetries. Multipliers
that fix the compression allow a further restriction on one designated
row, whose mask may be chosen least among its equivalent masks.

Within a unit, select one of the three ways of splitting its four row
lists into two pairs. Spectral bounds
reject pairs that cannot be completed by the opposite two lists. These
bounds use outward integer enclosures of the squared eigenvalues; the
error margin ensures that numerical rounding cannot exclude a solution.
The surviving pairs are joined using necessary autocorrelation and product
identities. Every matching candidate is tested by exact integer
arithmetic. Repeated classes and repeated rows are allowed throughout.

Finally, the absence of solutions in every unit, combined with the cover,
implies Theorem~\ref{thm:main}. This last step uses the correctness of the
executed search. The computation and its trust assumptions are described separately in
Section~\ref{sec:computation}.

\section{Rows, autocorrelations, and spectra}\label{sec:rows}
This section fixes the encoding and proves the necessary conditions used
to retain individual rows. All indices below are taken modulo the odd
order $n$, and $H=(n-1)/2$.

For example, at order $5$ a normalized symmetric first row has the form
$(1,a,b,b,a)$, with $a,b\in\{1,-1\}$. In general a normalized row is a
sequence $x=(x_i)_{i\in\Z/n\Z}$ with $x_0=1$ and $x_i=x_{-i}$.
Encode it by the $H$-bit integer $m_x$ whose bit $i-1$ is one exactly
when $x_i=-1$, for $1\leq i\leq H$. The function
$\operatorname{popcount}(m_x)$ counts the one bits of this mask.
Define its periodic autocorrelation by
\[
 \PAF_x(s)=\sum_{i\in\Z/n\Z}x_i x_{i+s}.
\]
For a symmetric circulant matrix with first row $x$, the entry in
position $(0,s)$ of its square is $\PAF_x(s)$. Thus the matrix equation is
exactly
\begin{equation}\label{eq:paf}
 \PAF_a(s)+\PAF_b(s)+\PAF_c(s)+\PAF_d(s)
 =\begin{cases}4n,&s=0,\\0,&s\ne0.\end{cases}
\end{equation}
Whenever a sum below runs over the four rows, it counts their four positions,
including repetitions. Autocorrelation is symmetric in $s$, so testing $0\leq s\leq H$ suffices.

\begin{lemma}\label{lem:rowsum}
If Williamson matrices of odd order $n$ exist, their first rows may be
normalized to $x_0=1$. Each normalized row has row sum
$r_x=n-4\operatorname{popcount}(m_x)$, and the four row sums satisfy
$r_a^2+r_b^2+r_c^2+r_d^2=4n$.
\end{lemma}
\begin{proof}
Negation preserves every autocorrelation. Negative entries in a normalized
symmetric row occur in pairs, giving the stated row sum. Summing
\eqref{eq:paf} over all shifts gives the sum of the squares of the row sums.
\end{proof}
At order $65$ the nondecreasing lists of absolute row sums are exactly
\[
 (1,3,5,15),\ (1,3,9,13),\ (3,3,11,11),\ (3,7,9,11),\ (7,7,9,9).
\]
Each sign is forced by $r_x\equiv65\pmod4$, giving the allowed signed
values $1,-3,5,-7,9,-11,13,-15$. Permuting the four rows therefore permits
sorting by the absolute row sums without loss.

For frequency $k\in\Z/n\Z$, the real eigenvalue of the circulant is
\[
 \lambda_x(k)=1+2\sum_{i=1}^{H}x_i\cos(2\pi ik/n).
\]
The Fourier transform diagonalizes all circulant matrices. Consequently
\eqref{eq:paf} is equivalent to
\begin{equation}\label{eq:spectral}
 \lambda_a(k)^2+\lambda_b(k)^2+\lambda_c(k)^2+\lambda_d(k)^2=4n
 \qquad(k\in\Z/n\Z).
\end{equation}
In particular every solution row satisfies $\lambda_x(k)^2\leq4n$.
The row generator considers every mask and retains only permitted row sums
and rows satisfying this necessary spectral bound, with the numerical
slack justified in Section~\ref{sec:arithmetic}.

\section{Compression and complete coverage}\label{sec:cover}
We now divide the possible solutions into finitely many units and justify
both symmetry reductions. These are reductions of the search domain;
they do not assert that a compressed solution necessarily lifts.

For $n=qL$, with compression factor $q$, define the compression of a row $x$ by
\[
 X_j=\sum_{t=0}^{q-1}x_{j+Lt},\qquad j\in\Z/L\Z.
\]
Here $q=5$ and $L=13$. A class consists of the retained rows with the
same signed row sum and compressed vector. Write
$\PAF_X(s)=\sum_{j\in\Z/L\Z}X_jX_{j+s}$ for compressed autocorrelation.

\begin{lemma}\label{lem:compression}
For four rows satisfying \eqref{eq:paf}, their compressions satisfy
\[
 \PAF_A(s)+\PAF_B(s)+\PAF_C(s)+\PAF_D(s)
 =\begin{cases}4n,&s=0\pmod L,\\0,&s\ne0\pmod L.\end{cases}
\]
\end{lemma}
\begin{proof}
Expanding the compressed autocorrelation gives
$\PAF_X(s)=\sum_{t=0}^{q-1}\PAF_x(s+Lt)$.
For a nonzero residue $s$ none of these shifts is zero modulo $n$.
For $s=0$ exactly one is zero. Summing over the four rows proves the claim.
\end{proof}
Call a tuple of four classes \emph{admissible} when its compressed
vectors satisfy Lemma~\ref{lem:compression} and its row sums satisfy
Lemma~\ref{lem:rowsum}. A \emph{lift} is a tuple of four retained rows,
one from each class; it need not itself solve the matrix equation.
Classes are ordered lexicographically by the seven entries
$(X_0,\ldots,X_6)$ of their symmetric compression. A class identifier is
its zero-based rank in this order; the data files use those identifiers. A \emph{sorted} class tuple
has nondecreasing absolute row sums and, within equal row sums,
nondecreasing class identifiers. Acting by a multiplier means applying it
to each compressed vector and then re-sorting within equal-row-sum
positions. The generator and cover check use this same order.

The admissible class tuples are enumerated by matching sums of two exact
compressed autocorrelation vectors to the complementary pair sums.
Every ordered class pair is included; equal classes are permitted.
Sorting positions of equal absolute row sum removes only row permutations.
The complete tuple set is also enumerated independently by matching the
exact integer vectors themselves.

A \emph{multiplier} is an invertible residue $u$ modulo $65$, acting on
a row by $x_i\mapsto x_{ui}$. It preserves \eqref{eq:paf}, since it
permutes the shifts. Its action on the compression depends only on its
residue modulo $13$. Opposite multipliers act identically on symmetric
rows. Every invertible residue modulo $13$ lifts to an invertible residue
modulo $65$: among $u,u+13,\ldots,u+52$, exactly one is divisible by $5$,
so the other four are coprime to $65$.

\begin{lemma}\label{lem:cover}
Let one representative of each orbit of admissible sorted class tuples
under $(\Z/13\Z)^*/\{1,-1\}$ be retained. Every order-$65$ Williamson
solution is equivalent, by a common multiplier and a permutation of its
rows, to a lift of a retained tuple.
\end{lemma}
\begin{proof}
A solution's tuple is admissible by Lemma~\ref{lem:compression}.
Choose a multiplier carrying it to its orbit representative, lift that
multiplier modulo $65$, and apply it to all four full rows. Re-sort rows
of equal absolute row sum. The resulting solution has the representative
tuple. All lifts of the representative are searched, even if more than one
multiplier carries the original tuple to it. Multipliers permute the exact
eigenvalues. Transformed solution rows therefore still satisfy
$\lambda_x(k)^2\leq260$ and pass the numerical filter with its proved
margin.
\end{proof}
A \emph{unit} is one representative in this lemma. The cover check computes
the exact compressed tuples and verifies membership of their representatives
in the supplied unit list. It checks invalid and duplicate representatives
as well as missing ones.

A further symmetry reduces one row list in each unit. The subgroup
\[
 K=\{u\in(\Z/65\Z)^*:u\equiv1\pmod{13}\}=\{1,14,27,53\}
\]
fixes every compressed vector. Choose one designated position in the unit
and retain only rows whose mask is least in its $K$-orbit.
Applying a common element of $K$ that minimizes that row preserves a
solution in the same unit. It is essential that this restriction applies
to only one designated row; independent minimization of all four rows
would not be justified.

\section{Filtering and matching full rows}\label{sec:join}
This section shows that every solution lifting a unit survives the pair
filters and is reached by the join. A pairing is one of the three splits
$(a,b)\mid(c,d)$, $(a,c)\mid(b,d)$, or $(a,d)\mid(b,c)$.
The implementation may choose any of these pairings and any designated row; these choices affect cost only.

\subsection{Conservative numerical enclosures}\label{sec:arithmetic}
We use integer enclosures to make the spectral comparisons safe under
rounding. At order $65$, the error in each computed squared eigenvalue is
below $4\cdot10^{-11}$ under the arithmetic assumptions and cosine-table
checks detailed in Appendix~\ref{app:numerics}. The row filter rejects only
when the computed square exceeds $260+10^{-6}$; it therefore retains
every solution row.

Let $\widetilde\lambda_x(k)$ be the computed eigenvalue. Define integer
lower and upper enclosures
\[
 \ell_x(k)=\max\{0,\lfloor64\widetilde\lambda_x(k)^2-0.01\rfloor\},
 \qquad
 h_x(k)=\lceil64\widetilde\lambda_x(k)^2+0.01\rceil.
\]
The error bound, including the two additional floating-point additions,
implies
\begin{equation}\label{eq:enclosure}
 \ell_x(k)\leq64\lambda_x(k)^2\leq h_x(k).
\end{equation}
Clamping the lower bound to zero preserves this enclosure, since a square
is nonnegative. No strictly positive gap is claimed after clamping.

\subsection{The pair filter}
We now use the opposite pair of row lists to bound the possible spectral
contribution of each pair. Consider two lists $\mathcal A,\mathcal B$ on one side and
$\mathcal C,\mathcal D$ on the other, after restricting the designated row.
Set $Q=64\cdot260$. For each frequency let
\[
 U(k)=\max\{0,Q-\min_{c\in\mathcal C}\ell_c(k)
                    -\min_{d\in\mathcal D}\ell_d(k)\},
\]
\[
 V(k)=\max\{0,Q-\max_{c\in\mathcal C}h_c(k)
                    -\max_{d\in\mathcal D}h_d(k)\}.
\]
Retain a pair $(a,b)$ only when, at every frequency,
\begin{equation}\label{eq:pair}
 \ell_a(k)+\ell_b(k)\leq U(k),\qquad
 h_a(k)+h_b(k)\geq V(k).
\end{equation}
The opposite side uses the same rule with the two sides interchanged.

\begin{lemma}\label{lem:pair}
Every pair belonging to a solution in the four lists passes
\eqref{eq:pair} on both sides.
\end{lemma}
\begin{proof}
For its complementary rows $c,d$, equation~\eqref{eq:spectral} gives
$64(\lambda_a(k)^2+\lambda_b(k)^2)=Q-64\lambda_c(k)^2-64\lambda_d(k)^2$.
Replacing each square by the appropriate lower or upper enclosure proves
the inequalities before clamping. Raising a negative upper bound to zero
only weakens the upper test. Raising a negative lower bound to zero changes
nothing, because $h_a+h_b$ is nonnegative.
\end{proof}
If a list is empty there is no solution in the unit. The implementation details are given in Appendix~\ref{app:join}.

\subsection{The product identity and the join}
The second matching condition is Williamson's classical product identity
\cite{Williamson,BKG}. We include its short argument to fix the normalization.

\begin{lemma}\label{lem:product}
For an odd-order Williamson solution normalized by $x_0=1$ in each row,
$a_i b_i c_i d_i=-1$ for every nonzero index $i$. Equivalently,
$m_a\mathbin{\oplus}m_b\mathbin{\oplus}m_c\mathbin{\oplus}m_d=2^H-1$,
where $\oplus$ is bitwise exclusive-or on $H$ bits.
\end{lemma}
\begin{proof}
In the integral group ring of the cyclic group of order $n$, write
$J$ for the sum of all group elements and $T_x$ for the sum of the
positions where row $x$ is negative. Then $x=J-2T_x$ and
$x^2=nJ-4|T_x|J+4T_x^2$. Symmetry and normalization make $|T_x|$ even.
Modulo $2$, the square $T_x^2$ equals the sum of the squares of its
individual group elements, because mixed terms occur twice. Reducing
$\sum_x x^2=4n\,1$ modulo $8$ and dividing by $4$ modulo $2$ gives
\[
 \sum_x\sum_{t\in T_x}t^2=1+J\pmod2.
\]
Squaring permutes a cyclic group of odd order. Every nonidentity position
therefore occurs in an odd number of the four negative sets, proving the
product identity and its mask form.
\end{proof}

For each row, form a linear hash of its nonzero-shift integer
autocorrelations,
\[
 f(x)=\sum_{s=1}^{H}R_s\PAF_x(s)\pmod {2^{64}},
\]
using fixed odd coefficients $R_s$; their reproducible choice is described
in Appendix~\ref{app:join}. A solution satisfies
$f(a)+f(b)=-f(c)-f(d)$ and
$m_a\oplus m_b=(2^H-1)\oplus m_c\oplus m_d$.
The search combines these two necessary equalities into a single
$64$-bit key, using the same function on both sides. Specifically it uses
$\alpha\oplus R\beta$ for hash sum $\alpha$ and mask value $\beta$, with a fixed odd $R$,
and maps a zero key to one. The right side uses the negative hash sum and
the $H$-bit complement of its mask XOR. Equal inputs give equal keys;
hash collisions may add candidates but cannot remove a solution.

The join retains every pair with an equal key, including duplicates.
Its storage and restart rules, justified in Appendix~\ref{app:join}, do
not remove any such pair.

Finally, each key match is checked using the four exact integer
autocorrelation vectors at shifts $0,\ldots,H$. Equation~\eqref{eq:paf}
then decides the matrix equation exactly. A separate witness checker can
also form the four matrices and multiply them directly. This last check
prevents false positives; it cannot repair a missing pair earlier in the
search.

\section{Computational result and reproducibility}\label{sec:computation}
This section states the finite outcome on which the theorem depends and
separates the checks of its records from re-execution of the search.
The programs and small data accompany the paper in \texttt{checks/} and
\texttt{data/}.

\subsection{The exhaustion record}
We state the finite computational result needed to finish the proof.
The CPU search regenerates its inputs from the problem definition. Its
counts are shown in Table~\ref{tab:counts}. A receipt is a record for one
unit: the four classes, pairing, designated row, list sizes, enumerated
pair counts, survivor counts, matching keys, solutions, and number of
passes. A pass processes one part of the pair lists, as defined in
Appendix~\ref{app:join}; every order-$65$ unit uses one pass. A key match
is a surviving left pair and a surviving right pair with equal keys.
A surviving pair is counted once for each side of each unit where it occurs.
The complete receipt collection has no unresolved unit.

\begin{table}[ht]
\centering
\begin{tabular}{lr}
\toprule
Quantity & Count\\
\midrule
Normalized masks examined & $2^{32}$\\
Retained rows & $396{,}050{,}997$\\
Compressed classes & $15{,}251$\\
Admissible sorted class tuples & $423{,}527$\\
Units after multiplier reduction & $70{,}699$\\
Enumerated pairs, both sides & $70{,}830{,}746{,}132{,}488$\\
Surviving pairs, both sides & $1{,}645{,}891{,}558{,}124$\\
Matches between pair keys & $0$\\
Solutions & $0$\\
\bottomrule
\end{tabular}
\caption{Counts in the complete CPU exhaustion record at order $65$.}
\label{tab:counts}
\end{table}

\begin{proof}[Computer-assisted conclusion of Theorem~\ref{thm:main}]
Every solution may be normalized and belongs to the retained row set by
Lemma~\ref{lem:rowsum} and the proved spectral margin.
Lemmas~\ref{lem:compression} and \ref{lem:cover} place an equivalent
solution in one of the covered units. The restriction on the designated
row preserves a solution in that unit. Lemmas~\ref{lem:pair} and
\ref{lem:product} show that its two pairs survive and have complementary mask values.
Moreover,
$ f(a)+f(b)+f(c)+f(d)
=\sum_{s=1}^{H}R_s(\PAF_a(s)+\PAF_b(s)+\PAF_c(s)+\PAF_d(s))=0$
modulo $2^{64}$, so their join keys are equal.
The exhaustive joins report no key match in any unit. Subject to the
search-execution trust base below, no solution exists.
\end{proof}

The complete CPU enumeration and two separately written complete GPU
searches all find no solution. They use the same mathematical compression
and symmetry reductions. The CPU implementation supplies its own row
generator, cosine calculations, class and unit enumeration, and pair join;
it does not read the GPU implementations' generated input files. The row
sets and class counts agree. Their pair filters differ, so the CPU and GPU
survivor counts are not expected to coincide and are not claimed to do so.
Appendix~\ref{app:gpu} documents the prior GPU computations and the
limits of their retained evidence. The per-unit comparisons below are
between the CPU search and the NumPy join, not between the GPU searches.
The public CPU implementation supplies a full replay without a GPU.

\subsection{Two independent record checkers}
The record checks establish the cover and its arithmetic accounting.
The first regenerates CPU inputs using the CPU generators, but checks the
cover with separately written code. The second has its own row, cosine,
cover, and record-checking code written from the mathematical specification.
Their row sets agree byte for byte, and both accept the complete cover
and pair totals. Neither rederives survivor or key-match counts.
Appendix~\ref{app:details} describes their scope and input sharing.

Here and throughout, independent means separately implemented, not
produced by separate human teams or independently replicated by outside
institutions. The AI provenance is stated on the first page.

Both checkers reject deliberately altered records and an incomplete
cover. Appendix~\ref{app:corruptions} states the changes tested.

These are \emph{coverage and accounting} checks. Neither checker proves
that every pair was correctly filtered, stored, and matched by the search.
In particular, a search bug could lose all genuine matches while leaving
consistent pair counts. The receipts are an exhaustion record, not a
short independently checkable proof of all the unsuccessful joins.

\subsection{Known cases and independent joins}
Known solutions and independently computed pair joins test that the
search can both recover solutions and reject false candidates. A
\emph{baseline enumeration} means a reference search with fewer
reductions. The baseline used here, without the search-time spectral, compression,
or product reductions, reproduces the published equivalence-class counts
at every odd order from $3$ through $33$ \cite[Table 1]{HKT}:
\[
\begin{array}{c|rrrrrrrr}
n&3&5&7&9&11&13&15&17\\ \hline
\#&1&1&2&3&1&4&4&4\\[2pt]
n&19&21&23&25&27&29&31&33\\ \hline
\#&6&7&1&10&6&1&2&5
\end{array}
\]
The equivalences used for odd-order class comparison are row negation,
row permutation, and a common invertible multiplier, as in
\cite[Section 2.2]{BKG}. The trivial order $1$ has one class as well.

Controls at orders $9$ through $35$ compare the reduced search with the
baseline, switching individual reductions off, removing both row sorting
and orbit reduction, and using compression alone. The corresponding solution sets agree after taking their
\emph{symmetry closures}: all normalized quadruples obtainable by
permuting rows and applying a common invertible multiplier. The order-$35$
reference uses a spectral row filter and yields no solution.

A \emph{row-filtered baseline} additionally applies the single-row
spectral restriction, but no pair filter or product-key restriction.
At order $39$ the controls recover the single published class and compare
with a row-filtered baseline. At order $45$ multiple compression factors
recover its single class; the variant comparisons use the symmetry closure
of the search output as reference, so those comparisons are consistency
checks, not an independent classification. At order $55$, both compression
factors $5$ and $11$ discover the single published class without receiving
a witness as input. In particular factor $5$ tests the same compression
factor as order $65$. These positive controls show that the search can
recover solutions, including repeated rows, where they exist.

A separately written implementation using NumPy arrays reconstructs row
lists, filters, and joins without sharing search code or using the product
identity. It compares the list sizes and survivor counts and performs its
own exact autocorrelation join. It agrees on all $3{,}309$ order-$45$ units
and recovers that order's solution.

At order $65$, it checks $300$ units sampled uniformly without replacement
with NumPy's \texttt{default\_rng(65)}, together with the $12$ units having
the largest enumerated-pair counts. These two sets are disjoint. The $312$
units contain $7{,}955{,}197{,}270$ surviving-pair instances, counted as above, about $0.4833\%$
of the total in Table~\ref{tab:counts}. Every sampled join gives zero
solutions, and all list sizes agree with the receipts.

Both survivor counts agree with the records on every checked unit.
The bookkeeping for numerical rounding boundaries, and further order-$55$
checks, is given in Appendix~\ref{app:details}.
These are sample-based implementation tests, not
a substitute for the complete enumeration. A positive matching test on
order-$65$ data, described in Appendix~\ref{app:positive}, also checks
that deliberately shortened keys produce the same nonzero match count
in separate join implementations.

\subsection{Public replay and data}
The public commands distinguish checking the record from repeating the
complete search. From the public folder, install \texttt{checks/requirements.txt} and invoke
\texttt{python3 checks/check.py}; the full requirements and mode descriptions
are in \texttt{checks/README.md}. The default check regenerates the order-$65$
rows and cover with both checkers and validates all receipts and corruption
controls. It also checks known-order counts, witness validity at orders $39$,
$45$, and $55$, and discovery at $39$ and $45$. It does not repeat all
order-$65$ pair joins. The \texttt{--full-controls} option adds the
unfiltered baseline counts through $33$ and discovery at $55$; the
\texttt{--full-search} option repeats every order-$65$ CPU join. The default mode checks the supplied records; the last mode recomputes
their mathematical results.

The full retained mask array alone has $1{,}584{,}203{,}988$ bytes,
one little-endian unsigned 32-bit integer per retained row. An optional dataset deposit could hold the complete row data, class and
unit data, full receipt collection, and control outputs for comparison. The compact class, unit, and receipt data needed by the default
check are already included in the public folder; the dataset is not required
for that check. The optional deposit contents and replay commands are specified in
\texttt{data/README.md}. The row data are regenerated from the public programs; no external
dataset is needed to run them. No separate dataset deposit exists and no
dataset DOI has been assigned. The complete CPU
search recorded about $271$ CPU-hours for its joins; this is a scale
estimate for a replay, not the cost of the short record checks.

\subsection{Trust boundary}
The theorem depends on the written reductions, the completeness of the
row and cover enumerations, and correct execution of the exhaustive
search. It also depends on the compiler, hardware, integer arithmetic,
floating-point model, and interval calculations used for the numerical
margins. Independent implementations, corruption rejection, and positive
controls supply evidence against implementation errors, but the search
has not been formally verified. No claim of a Lean proof or a short
standard unsatisfiability certificate is made. A reader seeking to repeat
the decisive computation must execute the full search, not only check
its receipts.

\appendix
\section{Numerical error bound}\label{app:numerics}
This appendix supplies the arithmetic details behind the conservative
spectral enclosures used in Section~\ref{sec:arithmetic}.

The row generator and the search evaluate the cosine table in IEEE binary64
arithmetic, rounded to nearest, without fast-math transformations.
The public build invokes GCC with \texttt{-O3 -fopenmp} and links
\texttt{libm}; it invokes G++ with \texttt{-O3 -std=c++17 -fopenmp}
for the second checker. The recorded CPU search additionally used
\texttt{-march=native}. No build uses \texttt{-ffast-math}. The numerical
bound assumes ordinary IEEE binary64 operations rather than x87 extended
intermediate precision; a build changing the arithmetic model requires
its own justification. The actual table entries are checked against interval evaluations of cosine
at $50$ decimal digits. Their errors are at most $2^{-49}$.
For order $65$ the eigenvalue sum is split into two sums of at most $16$
terms. Each partial sum has magnitude at most $16$, so the error in one
half is bounded by
\[
 16\,2^{-49}+16^2\,2^{-53}<6\cdot10^{-14}.
\]
Combining the halves, doubling, and adding one gives eigenvalue error below
$2.6\cdot10^{-13}$. Since $|\lambda_x(k)|\leq65$, the error in the
computed square, including rounding of the multiplication, is below
$4\cdot10^{-11}$. Fusing a multiplication and addition does not enlarge
this bound. The row filter rejects only when the computed square exceeds
$260+10^{-6}$, so it retains every solution row.


\section{Storage and key implementation}\label{app:join}
This appendix explains why the implemented key and data structures
preserve every pair equality required by a solution.
The coefficients $R_s$ are the successive outputs of the public
\texttt{splitmix64} function, initialized with the hexadecimal seed
\texttt{5713a1d2c0ffee65}, with the low bit set to one. The outputs are indexed from $s=0$ to $32$, with indices $1$ through $32$
used in $f$. The fixed multiplier $R$ is hexadecimal
\texttt{d6e8feb86659fd93}. The complete key function is specified in
\texttt{checks/search.c}. Their values affect collisions and cost, but
the argument needs only that both sides use the same function.

The pair bounds are formed in signed integer arithmetic before conversion
to unsigned 16-bit values. The stored row values are below $16{,}642$;
all intervening sums and differences fit the signed integer type.
Frequency reordering and vectorized comparisons do not change the
predicates. In a vectorized upper test, a row with $\ell_x>U$ is discarded
first, before forming $U-\ell_x$; the lower difference is
$\max\{0,V-h_x\}$. The remaining frequencies are tested with signed
integer sums. Empty lists contain no pairs and cannot yield a solution.

Surviving pairs on one side are stored by linear probing in a table whose
size exceeds twice the number of entries. No entry is deleted or
overwritten. A probe continues through every equal key until the first
empty slot, so repeated keys and repeated rows are retained. A capacity
overflow abandons the whole attempt and restarts in passes partitioned by
the high bits of this same key on both sides. Equal keys always belong to
the same pass. Every reported order-$65$ unit completed in one pass.


\section{Corruption controls}\label{app:corruptions}
This appendix records deliberate alterations that the independent
checkers must reject. They test whether the checks detect the particular
coverage and accounting errors for which they are designed.

Corruption controls exercise both implementations: deletion of a receipt,
a duplicate unit, an incorrect row-list size, an incorrect pair count,
keys truncated below their prescribed $64$ bits, a missing cover representative, and a
false solution must all be rejected. Additional row-data corruptions test
that the enumerated classes have not silently lost a row. The planted
cover gap is tested by the exhaustive cover computation itself, by re-enumerating the admissible tuples.


\section{Positive matching control at order 65}\label{app:positive}
This control tests the matching path on a full order-$65$ unit whose
normal $64$-bit keys have no matches. It truncates keys to $32$ bits to
create deliberate collisions, leaving the row lists and spectral filters
unchanged. This is a diagnostic setting, not the setting used to conclude
nonexistence.

For unit $5200$, the four list sizes are
$56{,}372$, $13{,}092$, $12{,}776$, and $2{,}830$.
The two sides have $6{,}611{,}810$ and $1{,}872{,}629$ surviving pairs.
The public search and a separately implemented join both find exactly
$2{,}918$ matching $32$-bit keys and zero exact solutions. These matches
are false collisions, not Williamson witnesses. Changing the expected
match count makes the control reject. The default public check includes
this test, alongside the unmodified $64$-bit receipt checks.

For scale only, if the pair keys were independent uniform random
$64$-bit values, the expected total number of equal keys would be
\[
 \frac{\sum_U S_{U,1}S_{U,2}}{2^{64}}
 =\frac{6{,}493{,}019{,}596{,}516{,}340{,}568}{2^{64}}
 \approx0.3519872976,
\]
where $S_{U,1}$ and $S_{U,2}$ are the survivor counts on the two sides
of unit $U$. At $32$ bits the same hypothetical model gives about
$2{,}883$ matches for unit $5200$, compared with the observed $2{,}918$.
The actual keys are deterministic and structured: no uniformity or
independence is proved, and these ratios are not probability bounds or
evidence of complete exhaustion. The control's evidential content is the
exact agreement of the two implemented joins on a nonzero match count.

\section{Computation by row-sum type}\label{app:types}
The following breakdown of the CPU receipts shows how the search cost
is distributed among the five possible absolute row-sum types. Survivors
count both sides of each pair join; CPU times are rounded.
\begin{center}
\begin{tabular}{lrrr}
\toprule
Absolute row sums & Units & Surviving pairs & CPU-hours\\
\midrule
$(1,3,5,15)$ & 19{,}793 & $439{,}857{,}850{,}463$ & 55.946 \\
$(1,3,9,13)$ & 21{,}291 & $528{,}816{,}451{,}429$ & 102.765 \\
$(3,3,11,11)$ & 4{,}733 & $103{,}766{,}383{,}699$ & 16.208 \\
$(3,7,9,11)$ & 19{,}692 & $449{,}073{,}745{,}221$ & 70.162 \\
$(7,7,9,9)$ & 5{,}190 & $124{,}377{,}127{,}312$ & 26.245 \\
\bottomrule
\end{tabular}
\end{center}

\section{Record-checking and control details}\label{app:details}
This appendix specifies implementation sharing and explains the numerical
comparison flags omitted from the main computational summary.

The record checks establish the cover and its arithmetic accounting. The
first checker reuses the CPU row and unit generators to regenerate their
outputs, constructs a separate
exact-vector cover enumeration, computes its own count of minimal masks
under $K$, and checks every receipt against these data. In particular,
each unit appears exactly once, its listed row counts agree, and each
enumerated pair count is the product of the corresponding list sizes.
The checker rejects records whose option fields do not state the full
$64$-bit key and prescribed necessary filters, and requires each recorded
solution count to be zero. It checks the recorded options; this alone
does not establish that the executable honored them. Survivor
and key-match counts are read from the records, not independently derived
by this checker. Its cover code does
not import the unit generator or the search.

The second checker was written from the mathematical specification
without using the first implementation's code. It shares the problem
definition and the supplied records, not the row generator, cosine table,
cover-enumeration code, or receipt-checking code. It enumerates all $2^{32}$ masks with its own C++ code and
integer cosine enclosures, obtains the same row set, performs its own
exact cover enumeration, and validates all unit records. Its regenerated row file agrees byte for byte with the first
implementation's. Both checkers accept the same cover and receipt totals.
They do not rederive the survivor or key-match counts.


Keys truncated to
$12$ bits test collisions. A constrained-capacity test exercises restarts;
its effective initial capacity is $1{,}024$ entries. 
When a value $64\widetilde\lambda_x(k)^2\pm0.01$ is within $10^{-7}$
of an integer, different summation orders may change its rounded enclosure.
The automatic survivor-count comparison is skipped for that unit, while
the independent solution search still runs. This affects $103$ of the
$312$ sampled units. Comparing the printed integer counts separately shows
that both survivor counts agree with their receipts in all $312$ units,
including those $103$.
At order $55$, the same implementation checks $300$ units sampled with
seed $55$ and the additional known-solution unit. It recovers the solution;
the printed survivor counts agree in all $301$ units, including the $29$
randomly sampled units whose automatic count comparison is skipped.

\section{Corroborating GPU computations}\label{app:gpu}
This appendix distinguishes the earlier complete GPU searches from the
CPU implementation distributed with this paper. Their results are
corroborating computational evidence, not formal verification of the CPU
program or per-unit agreement with its stronger pair filter.

Both GPU computations covered all $70{,}699$ compressed units and found
zero solutions. We call them the \emph{primary GPU search} and the
\emph{independent GPU replay}. The primary GPU search recorded about
$23.805$ GPU-hours of full-search time in its portable progress summary.
The independent GPU replay used a separately written implementation,
regenerated the row sets and checked the complete unit cover before enumerating
$146{,}369{,}334{,}587{,}912$ ordered pairs. It retained
$3{,}610{,}911{,}823{,}878$ pairs and found $30$ hash matches, all rejected
by exact autocorrelation checks. Its completed search intervals contain
each unit exactly once. This independent GPU replay recorded about
$2.060$ GPU-hours of search time and about $2.143$ occupied GPU-hours
including setup and controls. These resource measures are historical;
they are not controlled performance comparisons between algorithms.

The independent GPU replay used an integer lower bound for the pair
spectrum and did not use either the product-identity key or the
one-designated-row multiplier reduction. It therefore enumerated a
different number of pairs and retained more pairs than the CPU search.
Both full searches agree with the CPU search on the mathematical answer,
not on survivor or match totals. Their inputs, recorded build choices,
completed unit intervals, and outputs were checked against the retained
source and data records. Positive discovery controls at orders $25$,
$35$, and $45$ gave the known $10$, $0$, and $1$ equivalence classes.

The earlier GPU sources and output collections are not distributed in
this public CPU bundle. The portable record of the primary GPU search does not
recheck all original raw execution records, and neither CUDA program is
formally verified. The later CPU search is a separately coded, complete
repetition of the mathematical exhaustion; the fast public record checks
are not another complete repetition. All implementations share the
mathematical specification, and no independent human-team or institutional
replication is claimed. The self-contained reproduction route supplied
here is the complete CPU search, with the trust boundary stated in
Section~\ref{sec:computation}.

\begin{thebibliography}{99}
\bibitem{Williamson}
J. Williamson, Hadamard's determinant theorem and the sum of four squares,
Duke Mathematical Journal \textbf{11} (1944), 65--81.
\url{https://doi.org/10.1215/S0012-7094-44-01108-7}.
\bibitem{Turyn}
R. J. Turyn, An infinite class of Williamson matrices,
Journal of Combinatorial Theory, Series A \textbf{12} (1972), 319--321.
\url{https://doi.org/10.1016/0097-3165(72)90095-7}.
\bibitem{BGH}
L. Baumert, S. W. Golomb and M. Hall, Jr., Discovery of an Hadamard
matrix of order 92, Bulletin of the American Mathematical Society
\textbf{68} (1962), 237--238.
\url{https://doi.org/10.1090/S0002-9904-1962-10761-7}.
\bibitem{BH}
L. Baumert and M. Hall, Jr., Hadamard matrices of the Williamson type,
Mathematics of Computation \textbf{19} (1965), 442--447.
\url{https://doi.org/10.2307/2003674}.
\bibitem{Djokovic1992}
D. \v Z. \DJ okovi\'c, Williamson matrices of orders $4\cdot29$ and $4\cdot31$,
Journal of Combinatorial Theory, Series A \textbf{59} (1992), 309--311.
\url{https://doi.org/10.1016/0097-3165(92)90071-2}.
\bibitem{Djokovic1995}
D. \v Z. \DJ okovi\'c, Note on Williamson matrices of orders $25$ and $37$,
Journal of Combinatorial Mathematics and Combinatorial Computing
\textbf{18} (1995), 171--175.
\url{https://combinatorialpress.com/jcmcc/vol18/}.
\bibitem{Djokovic}
D. \v Z. \DJ okovi\'c, Williamson matrices of order $4n$ for $n=33,35,39$,
Discrete Mathematics \textbf{115} (1993), 267--271.
\url{https://doi.org/10.1016/0012-365X(93)90495-F}.
\bibitem{HKT}
W. H. Holzmann, H. Kharaghani and B. Tayfeh-Rezaie,
Williamson matrices up to order 59, Designs, Codes and Cryptography
\textbf{46} (2008), 343--352; a 2007 manuscript.
\url{https://doi.org/10.1007/s10623-007-9163-5}.
\bibitem{LS}
W. Lang and E. Schneider, Turyn type Williamson matrices up to order 99,
Designs, Codes and Cryptography \textbf{62} (2012), 79--84;
first published online March 5, 2011.
\url{https://doi.org/10.1007/s10623-011-9492-2}.
\bibitem{Zhang}
S. Y. Zhang, Programmable Emergence via Cognitive Engine: Longitudinal
Evidence of AI Research Maturation, April 2, 2026.
\url{https://doi.org/10.5281/zenodo.19392279}.
\bibitem{DK}
D. \v Z. \DJ okovi\'c and I. S. Kotsireas,
Compression of periodic complementary sequences and applications,
Designs, Codes and Cryptography \textbf{74} (2015), 365--377;
arXiv:1302.0571, first posted February 4, 2013.
\url{https://doi.org/10.1007/s10623-013-9862-z}.
\bibitem{BKG}
C. Bright, I. Kotsireas and V. Ganesh,
Applying Computer Algebra Systems with SAT Solvers to the Williamson
Conjecture, arXiv:1804.01172, first posted April 3, 2018;
Journal of Symbolic Computation \textbf{100} (2020), 187--209.
\url{https://arxiv.org/abs/1804.01172}.
\url{https://doi.org/10.1016/j.jsc.2019.07.024}.
\bibitem{BBCN}
A. Bennett, C. Bright, P. Colinot and A. Nayak,
Quaternionic Perfect Sequences and Hadamard Matrices,
arXiv:2601.22337, first posted January 29, 2026.
\url{https://arxiv.org/abs/2601.22337}.
\bibitem{BrightStatement}
C. Bright, Extended Research Statement, March 20, 2019.
\url{https://cs.curtisbright.com/extended-research-statement.pdf}.
\bibitem{EKLO}
R. Egan, I. Kotsireas, S. London and P. \'O Cath\'ain,
Recent progress in the construction of Hadamard matrices,
arXiv:2610.27175, first posted October 19, 2026.
\url{https://arxiv.org/abs/2610.27175}.
\bibitem{KMTZ}
H. Kharaghani, A. Mohammadian, B. Tayfeh-Rezaie and V. Zaitsev,
A Search for Hadamard Matrices of Near-Williamson Type,
arXiv:2605.08661v2, September 6, 2026.
First posted May 9, 2026 as version 1 by H. Kharaghani, A. Mohammadian
and B. Tayfeh-Rezaie, titled \emph{A search for Hadamard matrices of Williamson type}.
\url{https://arxiv.org/html/2605.08661v2}.
\end{thebibliography}
\end{document}
